\documentclass[12pt,a4paper]{mwrep}

\input{/home/tabaka/Documents/GitHub/Defs/defs.tex}

%\pagestyle{empty} 			%usuwa nr strony
\usepackage[polish]{babel}

\newtheorem*{remark}{Uwaga}
\title{Paradoks nieobecności}
\author{Arkadiusz Dąbal}
\date{\today}
	
\begin{document}
	\renewcommand{\proofname}{Dowód}
	\maketitle
	\thispagestyle{empty}
		
	\newpage
	
	\tableofcontents
		
	\chapter*{Wstęp}
	
	\newpage
	
	\chapter{Wiadomości wstępne}
	
	\Komentarz{Czy musi się znajdować tutaj 'rozdział'?}
	
	\section{Podstawowe pojęcia i oznaczenia}
	
	W tym rozdziale wprowadzamy podstawowe oznaczenia oraz pojęcia z zakresu teorii wyboru, które będą wykorzystywane w dalszej części pracy. Formalizacja ta umożliwia precyzyjne sformułowanie oraz analizę paradoksu nieobecności.
	
	Niech \(K\) oznacza zbiór kandydatów, przy czym \(\#K = m\), natomiast \(W\) oznacza zbiór wyborców, gdzie \(\#W = n\).
	
	\begin{definition}[elektorat]
		\label{def:elektorat}
		Zbiorem wszystkich elektoratów nazywamy zbiór
		\[
		\mathcal{E}(W) := \mathcal{P}(W) \setminus \{\emptyset\}.
		\]
		Elementy zbioru \(\mathcal{E}(W)\) interpretujemy jako niepuste podzbiory zbioru \(W\), czyli możliwe grupy wyborców uczestniczących w głosowaniu.
	\end{definition}
	
	\section{Preferencje jednostkowe}
	
	\begin{definition}[relacja preferencji]
		\label{def:relacja-preferencji}
		Relacją preferencji na zbiorze \(K\) nazywamy zupełny i przechodni porządek liniowy określony na \(K\).
	\end{definition}
	
	Relacja preferencji opisuje indywidualne upodobania wyborcy względem kandydatów. Zupełność oznacza, że każda para kandydatów jest porównywalna, natomiast przechodniość zapewnia spójność preferencji.
	
	\begin{definition}[relacja preferencji wyborcy]
		\label{def:relacja-preferencji-wyborcy}
		Dla każdego wyborcy \(w \in W\) jego relację preferencji oznaczamy przez \(\overset{w}{\succ}\). Zapis
		\[
		a \overset{w}{\succ} b
		\]
		oznacza, że wyborca \(w\) silnie preferuje kandydata \(a\) względem kandydata \(b\).
	\end{definition}
	
	\begin{definition}[przestrzeń relacji preferencji]
		\label{def:przestrzen-relacjo-preferencji}
		Przestrzenią relacji preferencji na zbiorze \(K\) nazywamy zbiór wszystkich relacji preferencji na \(K\), oznaczany przez
		\[
		\mathcal{R}.
		\]
	\end{definition}
	
	Dla każdego wyborcy \(w \in W\) jego relacja preferencji należy do przestrzeni \(\mathcal{R}\), tj.
	\[
	\overset{w}{\succ} \in \mathcal{R}.
	\]
	
	Dla zwięzłości zapisu pomijamy indeks wyborcy, gdy nie prowadzi to do niejednoznaczności. Wówczas ciąg symboli
	\[
	abcd
	\]
	oznacza relację preferencji postaci
	\[
	a \overset{w}{\succ} b \overset{w}{\succ} c \overset{w}{\succ} d.
	\]
	
	Równoważnie można ją przedstawić w postaci pionowej:
	\[
	\begin{array}{|c|}
		\hline
		w\\
		\hline
		a\\
		b\\
		c\\
		d\\
		\hline
	\end{array}
	\]
	
	\begin{definition}[słaba preferencja]
		\label{def:slaba-preferencja}
		Niech \(\overset{w}{\succ}\) będzie relacją preferencji wyborcy \(w \in W\).
		Relacją słabej preferencji związaną z \(\overset{w}{\succ}\) nazywamy relację \(\overset{w}{\succsim}\) określoną wzorem
		\[
		a \overset{w}{\succsim} b
		\iff
		a \overset{w}{\succ} b \ \text{lub}\ a \overset{w}{\sim} b,
		\]
		przy czym zapis \(a \overset{w}{\sim} b\) oznacza, że wyborca \(w\) jest obojętny pomiędzy kandydatami \(a\) i \(b\).
		Mówimy wtedy, że wyborca \(w\) słabo preferuje kandydata \(a\) nad kandydatem \(b\).
	\end{definition}
	[TODO - czy to jest potrzebne?]
	
	W dalszej części pracy rozważamy ustalony elektorat \(V \in \mathcal{E}(W)\).
	
	\section{Profile preferencji}
	
	\begin{definition}[profil preferencji]
		\label{def:profil-preferencji}
		Profilem preferencji na zbiorze wyborców \(W\) nazywamy funkcję
		\[
		R : W \to \mathcal{R},
		\]
		która każdemu wyborcy \(w \in W\) przyporządkowuje jego relację preferencji \(R(w)\).
	\end{definition}
	
	Zbiór wszystkich profili preferencji określonych na zbiorze \(W\) oznaczamy przez
	\[
	\mathcal{R}^{W}.
	\]
	
	Dla przejrzystości dalszych rozważań profile preferencji będziemy przedstawiać w postaci tabel opisujących liczebności grup wyborców posiadających dane relacje preferencji. Przykładowy profil przedstawiono w Tabeli~\ref{tab:przyklad_profilu}.
	
	\begin{table}[h]
		\centering
		\[
		\begin{array}{|c|c|}
			\hline
			1 & 3\\
			\hline
			a & d\\
			b & c\\
			c & b\\
			d & a\\
			\hline
		\end{array}
		\]
		\caption{Przykładowy profil preferencji}
		\label{tab:przyklad_profilu}
	\end{table}
	
	Pierwsza kolumna oznacza jednego wyborcę o preferencjach
	\[
	a \succ b \succ c \succ d,
	\]
	natomiast druga kolumna oznacza trzech wyborców o preferencjach
	\[
	d \succ c \succ b \succ a.
	\]
	
	\section{Model wyborczy}
	
	\begin{definition}[model wyborczy]
		\label{def:model-wyborczy}
		Modelem wyborczym nazywamy parę uporządkowaną
		\[
		\Sigma = (V, R|_V),
		\]
		przy czym \(V \in \mathcal{E}(W)\) jest elektoratem, natomiast \(R \in \mathcal{R}^{W}\) jest profilem preferencji na zbiorze wszystkich wyborców.
		
		Ograniczenie \(R|_V\) oznacza zawężenie profilu preferencji do wyborców należących do elektoratu \(V\).
	\end{definition}
	
	Zakładamy, że:
	\begin{itemize}
		\item ten sam profil preferencji prowadzi do tego samego wyniku wyborów,
		\item każdy dopuszczalny profil preferencji wyznacza pewien wynik wyborów.
	\end{itemize}
	
	\begin{definition}[metoda wyborcza]
		\label{def:metoda-wyborcza}
		Metodą wyborczą nazywamy funkcję
		\[
		M : \Sigma \to \mathcal{P}(K),
		\]
		która każdemu profilowi preferencji przyporządkowuje zbiór zwycięzców.
	\end{definition}
	
	\begin{definition}[metoda decyzyjna]
		\label{def:metoda-decyzyjna}
		Metodę wyborczą \(M\) nazywamy decyzyjną, jeżeli dla każdego profilu preferencji wskazuje dokładnie jednego zwycięzcę.
	\end{definition}
	
	\begin{definition}[metoda efektywna]
		\label{def:metoda-efektywna}
		Metodę wyborczą nazywamy efektywną, jeżeli dla każdego dopuszczalnego profilu preferencji wyłania co najmniej jednego zwycięzcę.
	\end{definition}
	
	\section{Operacje na elektoracie}
	
	W dalszej części pracy analizujemy wpływ udziału wyborców w głosowaniu na wynik wyborów. W tym celu wprowadzamy operacje modyfikacji elektoratu oraz odpowiadających im profili preferencji.
	
	Niech \(V \in \mathcal{E}(W)\) oraz \(R \in \mathcal{R}^{V}\). Zakładamy, że każdemu wyborcy \(w \in W\) przypisana jest relacja preferencji \(\overset{w}{\succ} \in \mathcal{R}\).
	
	Dla \(w \in V\) definiujemy usunięcie wyborcy:
	\[
	R - w := R \setminus \{(w, \overset{w}{\succ})\} \in \mathcal{R}^{V \setminus \{w\}}.
	\]
	
	Dla \(w \in W \setminus V\) definiujemy dodanie wyborcy:
	\[
	R + w := R \cup \{(w, \overset{w}{\succ})\} \in \mathcal{R}^{V \cup \{w\}}.
	\]
	
	Dla \(k \in \mathbb{N}_+\) zapisujemy skrótowo:
	\[
	R + k \cdot w,
	\qquad
	R - k \cdot w,
	\]
	gdzie oznacza to odpowiednio dodanie lub usunięcie \(k\) wyborców o relacji preferencji \(\overset{w}{\succ}\).
	
	Wprowadzone operacje pozwalają analizować zależność wyniku wyborów od składu elektoratu.
	
	
	
	
	
	
	
	
	
	
	
	%\begin{definition}[Metoda rozstrzygająca]
	%	Metodę wyborczą \(M\) nazywamy \textit{rozstrzygającą}, jeżeli każdemu profilowi preferencji \(R\) przyporządkowuje dokładnie jednego zwycięzcę, tj.
	%	\[
	%	M(R) \in K
	%	\]
	%	dla każdego dopuszczalnego profilu preferencji \(R\).
	%	
	%	Innymi słowy, metoda rozstrzygająca nie dopuszcza remisów ani zbioru współzwycięzców, lecz zawsze wskazuje jednego kandydata jako wynik wyborów.
	%\end{definition}
	
	\section{Kryterium Condorceta}
	
	\begin{definition}[kandydat Condorceta]
		\label{def:kandydat-condorceta}
		Niech \(R \in \mathcal{R}^{V}\) będzie profilem preferencji. Kandydat \(a \in K\) jest kandydatem Condorceta (zwycięzcą Condorceta), jeżeli w bezpośrednich porównaniach większościowych wygrywa z każdym innym kandydatem, tj.
		\[
		\forall b \in K,\ b \neq a \quad
		\#\{w \in V : a \overset{w}{\succ} b\}
		>
		\#\{w \in V : b \overset{w}{\succ} a\}.
		\]
	\end{definition}
	
	\begin{definition}[kryterium Condorceta]
		\label{def:kryterium-condorceta}
		Metoda wyborcza \(M\) spełnia kryterium Condorceta wtedy i tylko wtedy, gdy dla każdego profilu preferencji \(R\), jeśli istnieje kandydat Condorceta \(a \in K\), to jest on jedynym zwycięzcą wyłonionym przez metodę \(M\), tj.
		\[
		M(R) = \{a\}.
		\]
	\end{definition}
	
	\begin{definition}[przegrany Condorceta]
		\label{def:przegrany-condorceta}
		Niech \(R \in \mathcal{R}^{V}\). Kandydat \(a \in K\) jest przegranym Condorceta, jeżeli przegrywa z każdym innym kandydatem w porównaniach parami, tj.
		\[
		\forall b \in K,\ b \neq a \quad
		\#\{w \in V : b \overset{w}{\succ} a\}
		>
		\#\{w \in V : a \overset{w}{\succ} b\}.
		\]
	\end{definition}

	\begin{definition}[kryterium przegranego Condorceta]
		\label{def:kryterium-przegranego-condorceta}
		Metoda wyborcza \(M\) spełnia kryterium przegranego Condorceta wtedy i tylko wtedy, gdy dla każdego profilu preferencji \(R\), jeśli istnieje przegrany Condorceta \(a \in K\), to kandydat ten nie należy do zbioru zwycięzców wyłonionych przez metodę \(M\), tj.
		\[
		a \notin M(R).
		\]
	\end{definition}
	
	Kandydat będący przegranym Condorceta nie wygrywa żadnego porównania większościowego i stanowi naturalne przeciwieństwo zwycięzcy Condorceta.
	
	\begin{example}
		Rozważmy profil preferencji przedstawiony w Tabeli~\ref{tab:condorcet_przyklad}.
		
		\begin{table}[h]
			\centering
			\[
			\begin{array}{|c|c|c|}
				\hline
				2 & 3 & 2\\
				\hline
				a & b & c\\
				b & c & b\\
				c & a & a\\
				\hline
			\end{array}
			\]
			\caption{Profil preferencji dla trzech kandydatów}
			\label{tab:condorcet_przyklad}
		\end{table}
		
		W celu wyznaczenia relacji większościowych porównujemy kandydatów parami. Dla przykładu, w Tabeli~\ref{tab:porownanie_ab} wyróżniono pozycje kandydatów \(a\) oraz \(b\).
		
		\begin{table}[h]
			\centering
			\[
			\begin{array}{|c|c|c|}
				\hline
				2 & 3 & 2\\
				\hline
				\mathbf{a} & \mathbf{b} & c\\
				\mathbf{b} & c & \mathbf{b}\\
				c & \mathbf{a} & \mathbf{a}\\
				\hline
			\end{array}
			\]
			\caption{Porównanie większościowe kandydatów \(a\) i \(b\)}
			\label{tab:porownanie_ab}
		\end{table}
		
		Otrzymujemy następujące wyniki porównań większościowych:
		\[
		b \text{ pokonuje } a \text{ stosunkiem głosów } 5:2,
		\]
		\[
		a \text{ pokonuje } c \text{ stosunkiem głosów } 4:3,
		\]
		\[
		b \text{ pokonuje } c \text{ stosunkiem głosów } 5:2.
		\]
		
		Otrzymujemy więc relację większościową:
		\[
		b \succ a \succ c.
		\]
		
		Ponieważ kandydat \(b\) wygrywa z każdym innym kandydatem w bezpośrednich porównaniach większościowych, jest kandydatem Condorceta. Z kolei kandydat \(c\) przegrywa z każdym innym kandydatem, a więc jest przegranym Condorceta.
	\end{example}

	\begin{example}[paradox Condorceta]
		Rozważmy profil preferencji przedstawiony w Tabeli~\ref{tab:rps_condorcet}.
		
		\begin{table}[h]
			\centering
			\[
			\begin{array}{|c|c|c|}
				\hline
				2 & 3 & 2\\
				\hline
				a & b & c\\
				b & c & a\\
				c & a & b\\
				\hline
			\end{array}
			\]
			\caption{Profil preferencji typu kamień–papier–nożyce}
			\label{tab:rps_condorcet}
		\end{table}
		
		Otrzymujemy następujące porównania parami:
		\[
		a \succ b \quad (4:3),
		\]
		\[
		b \succ c \quad (5:2),
		\]
		\[
		c \succ a \quad (5:2).
		\]
		
		Zatem relacja większościowa ma cykliczną postać:
		\[
		a \succ b \succ c \succ a.
		\]
		
		W szczególności otrzymujemy błędne koło preferencji określane jako cykl condorcetowski \cite{rzazewski2014}.
	\end{example}

	\section{Kryterium maximina}
	
	\Komentarz{Dalej pozostał problem "maximina". Może faktycznie zapisze to jako maxi-mina.}
	
	\begin{definition}[margines większości]
		\label{def:magines-wiekszosci}
		Dla profilu preferencji \(R\) funkcją marginesu większości nazywamy funkcję
		\[
		g_R : K \times K \to \mathbb{Z},
		\]
		określoną wzorem
		\[
		g_R(a,b) =
		\#\{w \in W : a \overset{w}{\succ} b\}
		-
		\#\{w \in W : b \overset{w}{\succ} a\}.
		\]
	\end{definition}

	Funkcja \(g_R\) mierzy przewagę jednego kandydata nad drugim w porównaniach parami.
	
	\begin{lemma}[Własności funkcji marginesu większości]
		\label{lem:wlasnosci-marginesu}
		Dla dowolnego profilu preferencji \(R\) oraz kandydatów \(a,b,c \in K\) zachodzą następujące własności funkcji \(g_R\):
		
		\begin{enumerate}
			
			\item \textbf{Antysymetryczność:}
			\[
			g_R(a,b) = - g_R(b,a).
			\]
			
			\item \textbf{Ograniczoność wartości:}
			\[
			- |W| \leq g_R(a,b) \leq |W|.
			\]
			
			\item \textbf{Neutralność w porównaniu identycznych kandydatów:}
			\[
			g_R(a,a) = 0.
			\]
			
			\item \textbf{Zależność od relacji większościowej:}
			\[
			g_R(a,b) > 0 \iff a \text{ pokonuje } b \text{ w porównaniu większościowym}.
			\]
			
		\end{enumerate}
	\end{lemma}

	\begin{definition}[wartość minimalnego wyniku kandydata]
		\label{def:wartosc-minimalnego-wyniku-kandydata}
		Dla każdego kandydata \(k \in K\) definiujemy wartość
		\[
		m_k := \min_{x \in K \setminus \{k\}} g_R(k,x).
		\]
		Liczba \(m_k\) jest najsłabszym (najmniej korzystnym) wynikiem kandydata \(k\) w porównaniach parami.
	\end{definition}
	
	

	\begin{definition}[ważony turniej większościowy]
		Niech \(R\) będzie profilem preferencji na zbiorze kandydatów \(K\).
		Profil \(R\) indukuje ważony turniej skierowany \(T_R\), którego:
		\begin{itemize}
			\item wierzchołkami są kandydaci \(K\),
			\item każdej uporządkowanej parze \((x,y)\) przypisujemy krawędź skierowaną z wagą \(g_R(x,y)\).
		\end{itemize}
	\end{definition}

	\begin{example}[ważony turniej większościowy]
		\label{przyklad:truniej-wiekszosciowy}
		Rozważmy profil preferencji \(R\):
		\[
		\begin{array}{|c|c|c|c|}
			\hline
			2 & 1 & 1 & 1\\
			\hline
			a & b & c & d\\
			b & c & a & a\\
			c & a & b & b\\
			d & d & d & c\\
			\hline
		\end{array}
		\]
		
		Otrzymujemy następujące marginesy:
		\[
		g_R(a,b)=2,\quad g_R(a,c)=1,\quad g_R(a,d)=1,
		\]
		\[
		g_R(b,c)=0,\quad g_R(b,d)=1,\quad g_R(c,d)=1.
		\]
		
		Marginesy te możemy zapisać w postaci tabeli~\ref{tab:tabela_marginesow}.
		
		\begin{table}[h]
			\centering
			\[
			\begin{array}{c|cccc}
				g_R(x,y) & a & b & c & d\\
				\hline
				a & - & 2 & 1 & 1\\
				b & -2 & - & 0 & 1\\
				c & -1 & 0 & - & 1\\
				d & -1 & -1 & -1 & -
			\end{array}
			\]
			\caption{Macierz marginesów większości profilu \(R\)}
			\label{tab:tabela_marginesow}
		\end{table}
	
		Macierz z Tabeli~\ref{tab:tabela_marginesow} stanowi algebraiczną reprezentację ważonego turnieju, który indukowany jest przez profil \(R\). Indukowany ważony turniej \(T_R\) przedstawiono na Rysunku~\ref{fig:graf_turniej}.
		
		\begin{figure}[ht]
			\centering
			
			\begin{tikzpicture}[
				>=stealth,
				thick,
				vertex/.style={circle, draw, minimum size=1.2cm},
				edge label/.style={draw=none, fill=none, inner sep=1pt}
				]
				
				% Wierzchołki
				\node[vertex] (a) at (0,6) {$a$};
				\node[vertex] (b) at (6,6) {$b$};
				\node[vertex] (c) at (0,0) {$c$};
				\node[vertex] (d) at (6,0) {$d$};
				
				% krawędzie (majority relations)
				\draw[->] (a) -- node[edge label, above] {$2$} (b);
				\draw[->] (a) -- node[edge label, left] {$1$} (c);
				
				\draw[->] (a) -- node[edge label, sloped, above, pos=0.65] {$1$} (d);
				
				\draw[->] (b) -- node[edge label, right] {$1$} (d);
				\draw[->] (c) -- node[edge label, above] {$1$} (d);	
			\end{tikzpicture}
			\caption{Graf większościowy dla profilu \(R\)}
			\label{fig:graf_turniej}
		\end{figure}
	
		Strzałki w grafie reprezentują relacje większościowe między kandydatami: skierowanie krawędzi od \(x\) do \(y\) oznacza, że kandydat \(x\) pokonuje kandydata \(y\) w porównaniu parami. W przypadku braku strzałki między dwoma kandydatami mamy remis, tj. \(g_R(x,y)=0\).
		
		Z przedstawionego grafu można również bezpośrednio odczytać strukturę Condorcetowską profilu: kandydat \(a\) jest zwycięzcą Condorceta, natomiast kandydat \(d\) jest przegranym Condorceta.
	\end{example}
		
	\begin{definition}[zwyciężca maximina]
		Kandydata \(a \in K\) nazywamy zwycięzcą maximina, jeżeli maksymalizuje wartość
		\[
		\min_{k \in K \setminus \{a\}} g_R(a,k)
		\]
		przyjmuje wartość największą spośród wszystkich kandydatów z \(K\).
	\end{definition}

	Metodę która wybiera zwycięzców jako zbiór Kramera nazywamy metodą maximina.
	
	\begin{example}[zwycięzca Maximina]
		Rozważmy profil preferencji \(R\):
		\[
		\begin{array}{|c|c|c|c|}
			\hline
			2 & 3 & 3 & 1\\
			\hline
			a & b & c & d\\
			c & a & d & b\\
			b & c & b & c\\
			d & d & a & a\\
			\hline
		\end{array}
		\]
		
		Otrzymujemy macierz marginesów:
		\begin{table}[h]
			\centering
			\[
			\begin{array}{c|cccc|c}
				g_R(x,y) & a & b & c & d & m_x\\
				\hline
				a & - & \textbf{-5} & 1 & 1 & -5\\
				b & 5 & - & \textbf{-1} & 1 & -1\\
				c & \textbf{-1} & 1 & - & 7 & -1\\
				d & -1 & -1 & \textbf{-7} & - & -7
			\end{array}
			\]
			\caption{Macierz marginesów większości profilu \(R\)}
		\end{table}
		
		Zatem kandydaci \(b\) i \(c\) maksymalizują wartość minimalną wyniku i są zwyciężcami maximina.
	\end{example}
	
	\begin{wniosek}
		Każdy zwycięzca Condorceta jest zwycięzcą maximina.
	\end{wniosek}
	
	\begin{proof}
		Załóżmy, że \(a\) jest zwycięzcą Condorceta, tzn.
		\[
		g_R(a,x) > 0 \quad \text{dla każdego } x \neq a.
		\]
		Z antysymetryczności marginesu większości mamy
		\[
		g_R(x,a) = -g_R(a,x) < 0 \quad \text{dla każdego } x \neq a.
		\]
		Stąd dla każdego \(x \neq a\) istnieje kandydat \(a\), taki że
		\[
		g_R(x,a) < 0,
		\]
		a więc
		\[
		m_x = \min_{y \neq x} g_R(x,y) \le g_R(x,a) < 0.
		\]
		Zatem \(m_x < 0\) dla każdego \(x \neq a\).
		
		Ponadto
		\[
		m_a = \min_{y \neq a} g_R(a,y) > 0,
		\]
		co kończy dowód.
	\end{proof}
	
	W rozważanym wcześniej przykładzie~\ref{przyklad:truniej-wiekszosciowy}, kandydat \(a\) jest zwycięzcą Condorceta oraz spełnia on warunek
	
	\[
	m_a = 1 > 0,
	\]
	czyli kandydat \(a\) jest zwycięzcą maximina.
	
	Natomiast dla pozostałych kandydatów w tym przykładzie zachodzi
	\[
	m_b = -2, \qquad
	m_c = -1, \qquad
	m_d = -1,
	\]
	co pokazuje, że kandydaci niebędący zwycięzcą Condorceta mają ujemne wartości minimalnego wyniku.
	
	\begin{definition}[kryterium maximina]
		\label{def:kryterium_maximina}
		Metoda wyborcza \(M\) spełnia kryterium maximina, jeżeli dla każdego profilu \(R \in \mathcal{R}^W\), jeśli \(a\) jest zwyciężcą maximina, to \(M(R) = a\).
	\end{definition}

	\Komentarz{Metoda maximina \(\neq\) kryterium maximina. Musze to prawdopodobnie rozdzielić.
	
	Umieścić w dobrym miejscu w pracy definicje zbioru Kramera.}
	
	\begin{definition}[zbiór Kramera]
		Zbiór Kramera \(K^\ast\) definiujemy jako zbiór wszystkich kandydatów maksymalizujących wartość ich najgorszego wyniku w porównaniach parami, tj.
		\[
		K^\ast = \{a \in K : m_a \ge m_k \ \text{dla każdego } k \in K\}.
		\]
	\end{definition}
	
	Innymi słowy, \(K^\ast\) jest zbiorem kandydatów o najmniejszej maksymalnej porażce w porównaniach parami.
	
	\section{Kryterium Kemenego}
	
	\begin{definition}[ranking Kemeny’ego]
		\label{def:ranking-kemeny}
		Niech \(\mathcal R\) oznacza zbiór wszystkich liniowych porządków na zbiorze kandydatów \(K\).
		
		Ranking \(\succ \in \mathcal R\) nazywamy rankingiem Kemeny’ego dla profilu \(R\), jeżeli maksymalizuje wartość
		\[
		\sum_{w \in W}
		\#\{(a,b)\in K\times K \;:\; a \succ b \text{ oraz } a \succ_w b\}.
		\]
		
		Zbiór wszystkich rankingów Kemeny’ego dla profilu \(R\) oznaczamy przez \(\mathcal K(R)\).
	\end{definition}

	\Komentarz{Zapis poniżej \[
		\succ \in \mathcal K(R),
		\] jest fatalny i należy go poprawić.}
	
	\begin{definition}[zwycięzca Kemeny’ego]
		\label{def:zwyciezca-kemeny}
		Kandydata \(a \in K\) nazywamy zwycięzcą Kemeny’ego dla profilu \(R\), jeżeli istnieje ranking
		\[
		\succ \in \mathcal K(R),
		\]
		w którym kandydat \(a\) zajmuje pierwsze miejsce.
	\end{definition}

	\begin{example}
		\label{przyklad:kemeny}
		Rozważmy profil preferencji dany przez tabelę:
		
		\begin{table}[h]
			\centering
			\[
			\begin{array}{|c|c|c|}
				\hline
				3 & 2 & 2\\
				\hline
				a & b & c\\
				b & c & a\\
				c & a & b\\
				\hline
			\end{array}
			\]
			\caption{Profil preferencji dla metody Kemeny’ego}
			\label{tab:kemeny_example}
		\end{table}
	
		Otrzymujemy następujące porównania parami:
		\[
		a \succ b \quad (5:2),
		\]
		\[
		c \succ a \quad (4:3),
		\]
		\[
		b \succ c \quad (5:2).
		\]
		
		Rozważmy ranking
		\[
		a \succ b \succ c.
		\]
		
		Przeliczamy ile wyborców zgadza się z preferencjami
		\[a \succ b, \qquad b\succ c, \qquad a \succ c\],
		oraz wartości te sumujemy. Wówczas wartość tego porządku wynosi
		\[
		5 + 3 + 5 = 13,
		\]
		po przeliczeniu wszystkich porządków otrzymujemy:
		
		\[
		\begin{array}{c|c}
			\text{Ranking} & \text{Wartość Kemeny’ego}\\
			\hline
			a \succ b \succ c & 13\\
			a \succ c \succ b & 10\\
			b \succ a \succ c & 10\\
			b \succ c \succ a & 11\\
			c \succ a \succ b & 11\\
			c \succ b \succ a & 8
		\end{array}
		\]
		
		Ostatecznie porzadek \(a\succ b\succ c\) jest rankingiem Kemeny'ego, a kandydat \(a\) jest zwyciężcą Kemeny'ego.
	\end{example}

	\begin{definition}[kryterium Kemeny’ego]
		\label{def:kryterium_kemeny}
		Metoda wyborcza \(M\) spełnia kryterium Kemeny’ego, jeżeli dla każdego profilu \(R \in \mathcal{R}^W\), jeśli \(a\) jest zwycięzcą Kemeny’ego, to \(M(R) = a\).
	\end{definition}

	Metoda, która wybiera zwycięzców Kemeny'ego nazywamy metodą Kemeny'ego.

	\begin{wniosek}
		Każdy zwycięzca Condorceta jest zwycięzcą Kemeny’ego.
	\end{wniosek}
	
	\begin{proof}
		Załóżmy, że \(a\) jest zwycięzcą Condorceta dla profilu \(R\). Wówczas dla każdego kandydata \(x \neq a\) zachodzi
		\[
		a \succ x,
		\]
		tzn. większość wyborców preferuje \(a\) względem \(x\).
		
		Niech \(\succ\) będzie rankingiem Kemeny’ego. Załóżmy nie wprost, że kandydat \(a\) nie zajmuje w nim pierwszego miejsca, czyli istnieje taki kandydat \(b\), że
		\[
		b \succ a.
		\]
		Ponieważ jednak \(a\) wygrywa większościowo z każdym kandydatem, w szczególności z \(b\), zamiana pozycji kandydatów \(a\) oraz \(b\) zwiększa wartość funkcji Kemeny’ego. Otrzymujemy sprzeczność z maksymalnością wartością rankingu Kemeny'ego.
		
		Zatem istnieje ranking Kemeny’ego, w którym kandydat \(a\) zajmuje pierwsze miejsce, czyli \(a\) jest zwycięzcą Kemeny’ego.
	\end{proof}
	
	\chapter{Paradoks nieobecności}
	
	\section{Wprowadzenie}
	
	Paradoks nieobecności (ang.~\emph{no-show paradox}) jest jednym z istotnych i zarazem zaskakujących zjawisk badanych w teorii wyboru społecznego. Polega on na tym, że w przypadku pewnych metod wyborczych udział wyborcy w głosowaniu może prowadzić do wyniku mniej dla niego korzystnego niż jego absencja.
	
	Zjawisko to zostało po raz pierwszy opisane przez Fishburna i Bramsa \cite{fishburn1983} i od tego czasu stanowi ważny przedmiot badań w teorii systemów wyborczych. Co szczególnie istotne, Moulin \cite{moulin1988} wykazał, że szeroka klasa metod spełniających kryterium Condorceta jest podatna na ten paradoks, co często interpretuje się jako argument przeciwko stosowaniu tzw.~Condorcet extensions (czyli sprawdzeniu czy istnieje kandydat Condorceta i jeśli istnieje to wybraniu go jako zwyciężcę).
	
	W literaturze podkreśla się, że paradoksy wyborcze stanowią nieuniknioną konsekwencję ogólności modeli głosowania oraz braku możliwości jednoczesnego spełnienia wszystkich naturalnych postulatów racjonalności procedur wyborczych. Jednocześnie pozostaje otwarte pytanie, na ile zjawiska te mają znaczenie praktyczne.
	
	W szczególności nie jest jasne, jak często paradoks nieobecności występuje w typowych modelach preferencji oraz jak jego częstość zależy od liczby kandydatów i wyborców. Problem ten został podjęty w nowszych badaniach empiryczno-analitycznych, m.in. przez Brandta, Hofbauera i Strobla \cite{brandt2017nsp}, którzy analizują częstość występowania paradoksu nieobecności dla wybranych metod Condorcetowskich, wykorzystując zarówno metody analityczne oparte na teorii Ehrharta, jak i symulacje komputerowe.
	
	W niniejszej pracy wprowadzamy formalny model pozwalający na analizę wpływu obecności wyborcy na wynik wyborów oraz systematyczne badanie warunków występowania paradoksu nieobecności w różnych procedurach wyborczych.
	
	\section{Przykład paradoksu nieobecności}
	
	\begin{example}[paradoks nieobecności dla procedury głosowań parami]
		Rozważmy procedurę głosowań parami ze statusem quo równym kandydatowi \(d\), to znaczy, że jeśli pewna ilość wborców nie zagłosuje to zwyciężcą zostaje kandydat \(d\).
		Agenda głosowań ma postać:
		\[
		b \text{ vs } d,\qquad
		\text{następnie zwycięzca vs } a,\qquad
		\text{następnie zwycięzca vs } c.
		\]
		
		Niech profil preferencji będzie dany tabelą:
		\[
		\begin{array}{|c|c|c|}
			\hline
			10 & 10 & 10\\
			\hline
			a & b & d\\
			b & d & c\\
			d & c & a\\
			c & a & b\\
			\hline
		\end{array}
		\]
		
		\begin{itemize}
			\item W pierwszym głosowaniu porównywani są kandydaci \(b\) oraz \(d\). Kandydat \(b\) wygrywa stosunkiem głosów \(20:10\).
			
			\item Następnie kandydat \(b\) zostaje porównany z kandydatem \(a\). W tym głosowaniu zwycięża \(a\) stosunkiem głosów \(20:10\).
			
			\item W ostatnim etapie kandydat \(a\) zostaje zestawiony z kandydatem \(c\). Kandydat \(c\) wygrywa stosunkiem głosów \(20:10\), a zatem końcowym zwycięzcą procedury zostaje \(c\).
			
			\item Każda z trzech grup wyborców preferuje kandydata \(d\) względem zwycięzcy \(c\).
			
			\item Oznacza to, że wszyscy wyborcy preferowaliby utrzymanie status quo względem wyniku uzyskanego przy pełnej frekwencji.
			
			\item W konsekwencji pełny udział wyborców w procedurze prowadzi do wyniku gorszego z punktu widzenia wszystkich grup, co stanowi przykład paradoksu nieobecności.
		\end{itemize}
	\end{example}
	
	W tej procedurze kolejność głosowań ma istotny wpływ na wynik, przez co metoda nie jest neutralna.
	
	\section{Odporność na nieobecność}
	
	Intuicyjnie, paradoks nieobecności występuje wtedy, gdy obecność wyborcy w elektoracie prowadzi do wyniku wyborów, który jest dla niego mniej preferowany niż wynik uzyskany w sytuacji jego nieobecności.
	
	Oznacza to, że udział w głosowaniu może działać przeciwko interesowi samego głosującego, mimo że jego preferencje są w pełni uwzględniane w procesie decyzyjnym.
	
	\medskip
	
	Poniższy przykład ilustruje występowanie tego zjawiska dla metody Copelanda.
	
	\begin{example}[paradoks nieobecności dla metody Copelanda]
		Rozważmy profil preferencji:
		\[
		\begin{array}{|c|c|c|c|}
			\hline
			1 & 1 & 1 & 1\\
			\hline
			A & B & D & C\\
			C & A & A & D\\
			D & C & C & B\\
			B & D & B & A\\
			\hline
		\end{array}
		\]
		
		Załóżmy dodatkowo, że w przypadku remisu metoda wybiera kandydata \(C\).
		
		Rozważmy dodanie wyborcy o preferencjach:
		\[
		C \succ A \succ B \succ D.
		\]
		
		Otrzymujemy profil:
		\[
		\begin{array}{|c|c|c|c|c|}
			\hline
			1 & 1 & 1 & 1 & 1\\
			\hline
			A & B & D & C & C\\
			C & A & A & D & A\\
			D & C & C & B & B\\
			B & D & B & A & D\\
			\hline
		\end{array}
		\]
		
		W przypadku początkowego profilu kandydaci \(A\) oraz \(C\) uzyskują równą wartość funkcji Copelanda, a zgodnie z przyjętą regułą rozstrzygania remisów zwycięzcą zostaje \(C\).
		
		Po dodaniu nowego wyborcy jedynym zwycięzcą metody Copelanda zostaje kandydat \(A\).
		
		Nowy wyborca preferuje \(C \succ A\), a więc woli wynik uzyskany w sytuacji swojej nieobecności.
		
		Oznacza to, że jego udział w głosowaniu prowadzi do wyniku mniej preferowanego z jego punktu widzenia.
	\end{example}
	
	\medskip
	
	\begin{definition}[pozytywne kryterium obecności]
		Metoda wyborcza \(M\) spełnia pozytywne kryterium obecności, jeżeli dla każdego profilu preferencji \(R \in \mathcal{R}^{V}\) oraz każdego wyborcy \(w \in V\) zachodzi:
		
		\[
		\text{jeżeli } a \in M(R - w) \text{ oraz } a \overset{w}{\succsim} b \ \text{dla każdego } b \in K,
		\text{ to } a \in M(R).
		\]
		
		Innymi słowy, jeżeli kandydat najwyżej oceniany przez wyborcę \(w\) zostaje wybrany przy jego nieobecności, to pozostaje on zwycięzcą również przy jego udziale.
	\end{definition}

	\begin{definition}[negatywne kryterium obecności]
		Metoda wyborcza \(M\) spełnia negatywne kryterium obecności, jeżeli dla każdego profilu preferencji \(R \in \mathcal{R}^{V}\) oraz każdego wyborcy \(w \in V\) zachodzi:
		
		\[
		\text{jeżeli } a \in M(R) \text{ oraz } b \overset{w}{\succsim} a \ \text{dla każdego } b \in K,
		\text{ to } a \in M(R - w).
		\]
		
		Innymi słowy, jeżeli kandydat najniżej oceniany przez wyborcę \(w\) zostaje wybrany przy jego udziale, to pozostaje on zwycięzcą również przy jego nieobecności.
	\end{definition}
	
	Formalnie własność tę ujmujemy następująco.
	
	\begin{definition}[metoda odporna na nieobecność]
		Metodę wyborczą \(M\) nazywamy odporną na nieobecność, jeżeli dla każdego profilu preferencji \(R \in \mathcal{R}^{V}\) oraz każdego wyborcy \(w \in V\) zachodzi
		\[
		M(R) \overset{w}{\succsim} M(R - w).
		\]
	\end{definition}
	
	Paradoks nieobecności pozostaje w ścisłym związku z kryterium Condorceta, co zostanie wykorzystane w dalszej części pracy.
	
	\section{Twierdzenie Moulina}

	\begin{lemma}
		\label{lem:condorcet-minimax}
		Jeżeli kandydat \(a \in K\) jest kandydatem Condorceta, to \(a\) jest jedynym zwycięzcą metody maximina.
	\end{lemma}
	
	\begin{proof}
		Niech \(a \in K\) będzie kandydatem Condorceta, tzn.
		\[
		g_R(a,k) > 0 \qquad \text{dla każdego } k \in K \setminus \{a\}.
		\]
		
		Rozważmy wartości
		\[
		m_k := \min_{x \in K \setminus \{k\}} g_R(k,x),
		\qquad k \in K.
		\]
		
		Dla kandydata \(a\) mamy:
		\[
		m_a = \min_{k \neq a} g_R(a,k) > 0,
		\]
		ponieważ wszystkie wartości \(g_R(a,k)\) są dodatnie.
		
		Niech teraz \(k \neq a\) będzie dowolnym innym kandydatem. W szczególności zachodzi
		\[
		g_R(k,a) < 0,
		\]
		co implikuje
		\[
		m_k = \min_{x \neq k} g_R(k,x) \le g_R(k,a) < 0.
		\]
		
		Zatem dla każdego \(k \neq a\) mamy ścisłą nierówność
		\[
		m_a > m_k.
		\]
		
		W konsekwencji kandydat \(a\) osiąga największą wartość funkcji \(m_k\), czyli
		\[
		m_a = \max_{k \in K} m_k.
		\]
		
		Co więcej, ponieważ dla każdego \(k \neq a\) zachodzi \(m_k < m_a\), kandydat \(a\) jest jedynym maksymalizatorem tej wartości. A zatem \(a\) jest jedynym zwycięzcą metody maximina.
	\end{proof}

	\begin{theorem}[o istnienie metody Condorceta zgodnej z partycypacją dla \(\leq 3\) kandydatów]
		Jeżeli liczba kandydatów nie przekracza trzech, to istnieje metoda wyborcza spełniająca kryterium Condorceta oraz jednocześnie pozytywne i negatywne kryterium obecności.
	\end{theorem}
	
	\begin{proof}
		Rozważamy metodę wyborczą wybierającą zawsze kandydata ze zbioru Kramera, zgodnie z podejściem Moulin \cite{moulin1988}. Z Lematu~\ref{lem:condorcet-minimax} wynika, że metoda \(M\) spełnia kryterium Condorceta.
		
		Pokażemy teraz, że metoda \(M\) spełnia również pozytywne oraz negatywne kryterium obecności dla trzech kandydatów.
		
		Rozważmy dowolny profil preferencji oraz oznaczmy kandydatów przez \(a\), \(b\) oraz \(c\). Załóżmy, że metoda \(M\) wybiera kandydata \(a\). Oznacza to, że \(a \in K^\ast\), a więc w szczególności
		\[
		m_a \geq m_b \quad \text{oraz} \quad m_a \geq m_c.
		\]
		
		Rozważamy teraz zmianę profilu polegającą na dodaniu jednego wyborcy. Każda wartość \(m_k\) może zmienić się co najwyżej o \(1\), ponieważ pojedynczy głos wpływa jedynie o jednostkową zmianę na marginesy porównań parami.
		
		W związku z tym analizujemy wszystkie możliwe relacje preferencyjne nowego wyborcy względem kandydatów \(a,b,c\). Rozróżniamy trzy istotne przypadki.
		
		\begin{enumerate}
			\item Nowy wyborca ściśle preferuje \(a\) względem \(b\) oraz \(c\).
			
			\item Nowy wyborca słabo preferuje \(b\) względem \(a\), ale ściśle preferuje \(a\) względem \(c\).
			
			\item Nowy wyborca słabo preferuje zarówno \(b\), jak i \(c\) względem \(a\).
		\end{enumerate}
	
		\begin{enumerate}[\text{Ad} 1.]
			\item W pierwszym przypadku wartość \(m_a\) wzrasta o \(1\). Ponieważ zmiany wszystkich \(m_k\) są ograniczone jednostkowo, wartości \(m_b\) oraz \(m_c\) nie mogą wzrosnąć bardziej niż \(m_a\). W konsekwencji nadal zachodzi:
			\[
			m_a \geq m_b \quad \text{oraz} \quad m_a \geq m_c.
			\]
			
			Zatem kandydat \(a\) pozostaje w zbiorze Kramera po dodaniu wyborcy, a więc metoda \(M\) nadal go wybiera. Oznacza to spełnienie zarówno pozytywnego, jak i negatywnego kryterium obecności.
			
			\item W drugim przypadku wartość \(m_c\) maleje o \(1\), natomiast \(m_a\) może zmaleć co najwyżej o \(1\), więc nadal zachodzi:
			\[
			m_a \geq m_c.
			\]
			
			Oznacza to, że kandydat $c$ nie może wyprzedzić kandydata $a$ w rankingu wartości $m_k$, a więc nie może wejść do zbioru Kramera, jeśli wcześniej do niego nie należał.
			
			W szczególności dodanie nowego wyborcy nie prowadzi do wykluczenia kandydata $a$ z $K^\ast$, a metoda pozostaje zgodna zarówno z pozytywnym, jak i negatywnym kryterium obecności.
			
			\item W trzecim przypadku zmiany nie naruszają relacji preferencyjnych istotnych dla zbioru Kramera w sposób wymuszający wykluczenie kandydata \(a\), więc oba kryteria mogą być spełnione bez sprzeczności.
		\end{enumerate}

		Dla dwóch kandydatów własność ta wynika bezpośrednio z większościowego charakteru metody.
	\end{proof}

	\begin{lemma}\label{lem:condorcet_obecnosc}
		Załóżmy, że metoda \(M\) spełnia kryterium Condorceta oraz pozytywne kryterium obecności. Wówczas dla każdych różnych kandydatów \(a,b \in K\), każdego elektoratu \(W\) oraz każdego profilu preferencji \(\succ \in R^W\),
		\[
		m_a < 0
		\qquad \text{oraz} \qquad
		m_a - 2 \geq g_R(b,a)
		\]
		implikuje
		\[
		M(R^{W}) \neq b.
		\]
	\end{lemma}
	
	\begin{proof}
		Rozważmy dowolny elektorat \(W\), profil preferencji \(\succ \in R^W\) oraz kandydatów \(a,b \in K\), dla których
		\[
		m_a < 0
		\qquad \text{oraz} \qquad
		m_a - 2 \geq g_R(b,a),
		\]
		i załóżmy przeciwnie, że
		\[
		M(R^{W}) = b.
		\]
		
		Ponieważ \(m_a < 0\), kandydat \(a\) przegrywa co najmniej jedno porównanie parami, a zatem nie jest zwycięzcą Condorceta.
		
		Niech \(R^{V}\) będzie profilem otrzymanym z \(R^W\) przez dodanie \(1-m_a\) nowych wyborców, z których każdy umieszcza kandydata \(b\) na pierwszym miejscu, a kandydata \(a\) na drugim miejscu.
		
		\[
		\begin{array}{|c|}\hline
			1-m_a\\\hline
			b\\
			a\\
			\vdots
			\\\hline
		\end{array}
		\]
		
		W profilu \(R^W\) najgorszy wynik kandydata \(a\) ma wartość \(m_a\). Ponadto z założenia
		\[
		m_a - 2 \geq g_R(b,a),
		\]
		a więc
		\[
		g_R(a,b) = -g_R(b,a) \geq 2 - m_a.
		\]
		
		Dodanie \(1-m_a\) nowych wyborców zwiększa o \(1-m_a\) wynik kandydata \(a\) w porównaniach parami ze wszystkimi kandydatami różnymi od \(b\). Natomiast wynik porównania parami pomiędzy \(a\) i \(b\) zmniejsza się o \(1-m_a\).
		
		W szczególności wszystkie porażki parami kandydata \(a\) zostają wyeliminowane. Ponadto kandydat \(a\) nadal pokonuje kandydata \(b\), gdyż
		\[
		g_R(a,b) - (1-m_a) \geq 1.
		\]
		
		Zatem kandydat \(a\) staje się zwycięzcą Condorceta w profilu \(R^{V}\).
		
		Ponieważ jednak metoda \(M\) wybrała kandydata \(b\) w profilu \(R^W\), pozytywne kryterium obecności wymaga, aby kandydat \(b\) został wybrany również w profilu \(R^{V}\).
		
		Otrzymujemy sprzeczność z kryterium Condorceta, gdyż w profilu \(R^{V}\) zwycięzcą Condorceta jest kandydat \(a\). Kończy to dowód.
	\end{proof}
	
	\begin{theorem}
		Jeżeli liczba kandydatów spełnia \(\#K \geq 4\), to żadna metoda wyborcza nie spełnia jednocześnie kryterium Condorceta oraz pozytywnego kryterium obecności.
	\end{theorem}
	
	\begin{proof}
		Niech \(M\) będzie metodą spełniającą kryterium Condorceta oraz pozytywne kryterium obecności. Rozważmy czterech kandydatów:
		\[
		a,b,c,d.
		\]
		
		\begin{table}[h]
			\centering
			\[
			\begin{array}{|c|c|c|c|}
				\hline
				6 & 3 & 8 & 7\\
				\hline
				a & a & d & b\\
				d & d & b & c\\
				c & b & c & a\\
				b & c & a & d\\
				\hline
			\end{array}
			\]
			\caption{Pierwszy profil preferencji}
			\label{tab:pierwszy_profil}
		\end{table}
		
		Pierwszy profil przedstawiono w Tabeli~\ref{tab:pierwszy_profil}. Liczby nad kolumnami oznaczają liczbę wyborców posiadających ranking zapisany poniżej. Wszyscy pozostali kandydaci (jeżeli istnieją) są umieszczani poniżej kandydatów \(a,b,c,d\).
		
		\begin{figure}[h]
			\centering
			
			\begin{tikzpicture}[
				>=stealth,
				thick,
				vertex/.style={circle, draw, minimum size=1.2cm},
				edge label/.style={draw=none, fill=none, inner sep=1pt}
				]
				
				% Wierzchołki
				\node[vertex] (a) at (0,6) {$a$};
				\node[vertex] (b) at (6,6) {$b$};
				\node[vertex] (c) at (0,0) {$c$};
				\node[vertex] (d) at (6,0) {$d$};
				
				% krawędzie (majority relations)
				\draw[->] (b) -- node[edge label, above] {$6$} (a);
				\draw[->] (c) -- node[edge label, left] {$6$} (a);
				
				\draw[->] (a) -- node[edge label, sloped, above, pos=0.65] {$8$} (d);
				
				\draw[->] (d) -- node[edge label, right] {$10$} (b);
				\draw[->] (d) -- node[edge label, below] {$10$} (c);
				
				\draw[->] (b) -- node[edge label, sloped, below, pos=0.65] {$12$} (c);
				
			\end{tikzpicture}
		
		\end{figure}
		
		\begin{figure}[h]
			\centering
			\begin{tikzpicture}[
				>=stealth,
				thick,
				vertex/.style={circle, draw, minimum size=1.2cm}
				]
				
				% Wierzchołki
				\node[vertex] (a) at (0,6) {$a$};
				\node[vertex] (b) at (6,6) {$b$};
				\node[vertex] (c) at (0,0) {$c$};
				\node[vertex] (d) at (6,0) {$d$};
				
				% krawędzie
				\draw[->] (b) -- (a);
				\draw[->] (c) -- (a);
				
				\draw[->] (a) -- (d);
				\draw[->] (d) -- (b);
				\draw[->] (d) -- (c);
				\draw[->] (b) -- (c);
				
				% etykiety (OSOBNO – pełna kontrola)
				\node at (3,6.3) {$6$};   % b -> a
				\node at (-0.3,3) {$6$};  % c -> a
				
				\node at (5.2,1.8) {$8$};   % a -> d (przekątna)
				\node at (0.8,1.8) {$12$};  % b -> c (przekątna)
				
				\node at (6.3,3) {$10$};  % d -> b
				\node at (3,-0.3) {$10$}; % d -> c
				
			\end{tikzpicture}
		\caption{Graf większościowy dla profilu z Tabeli~\ref{tab:pierwszy_profil}}
		\label{fig:graf_pierwszy_profil}
		\end{figure}
		
		Rysunek~\ref{fig:graf_pierwszy_profil} przedstawia skierowany graf ważony opisujący marginesy zwycięstw i porażek parami pomiędzy czterema głównymi kandydatami. Krawędź skierowana z \(b\) do \(a\) o wadze \(6\) oznacza, że kandydat \(b\) pokonuje kandydata \(a\) przewagą \(6\) głosów.
		
		W profilu uczestniczy \(24\) wyborców. Mamy:
		\[
		m_a = -6
		\qquad \text{oraz} \qquad
		g_R(a,d)=8.
		\]
		Na mocy Lematu~\ref{lem:condorcet_obecnosc} metoda \(M\) nie może wybrać kandydata \(d\).
		
		Ponadto:
		\[
		m_d = -8
		\qquad \text{oraz} \qquad
		g_R(d,b)=10,
		\]
		więc metoda \(M\) nie może wybrać kandydata \(b\).
		
		Dla każdego kandydata
		\[
		x \in K \setminus \{a,b,c,d\}
		\]
		zachodzi:
		\[
		m_a=-6
		\qquad \text{oraz} \qquad
		g_R(a,x)=24,
		\]
		a więc metoda \(M\) nie może wybrać kandydata \(x\).
		
		Zatem zbiór zwycięzców ogranicza się do kandydata \(a\) albo \(c\).
		
		Dodajmy teraz ośmiu wyborców, otrzymując drugi profil preferencji. Każdy z nowych wyborców traktuje kandydatów \(a\) i \(c\) jako równorzędnych i umieszcza ich wspólnie na pierwszym miejscu. Następnie preferowany jest kandydat \(b\), potem kandydat \(d\), a wszyscy pozostali kandydaci znajdują się poniżej \(d\).
		
		\begin{table}[h]
			\centering
			\[
			\begin{array}{|c|c|c|c|c|}
				\hline
				6 & 3 & 8 & 7 & 8\\
				\hline
				a & a & d & b & a \\
				d & d & b & c & c \\
				c & b & c & a & b \\
				b & c & a & d & d \\
				\hline
			\end{array}
			\]
			\caption{Drugi profil preferencji (z dodanymi ośmioma wyborcami)}
			\label{tab:drugi_profil}
		\end{table}
		
		Odpowiadający temu profilowi graf przedstawiono na Rysunku~2.
		
		\begin{figure}[h]
			\centering
			
			\begin{tikzpicture}[
				>=stealth,
				thick,
				vertex/.style={circle, draw, minimum size=1.2cm},
				edge label/.style={draw=none, fill=none, inner sep=1pt}
				]
				
				% Wierzchołki
				\node[vertex] (a) at (0,6) {$a$};
				\node[vertex] (b) at (6,6) {$b$};
				\node[vertex] (c) at (0,0) {$c$};
				\node[vertex] (d) at (6,0) {$d$};
				
				% --- relacje większościowe (bazowe) ---
				\draw[->] (a) -- node[edge label, above] {$2$} (b);
				\draw[->] (c) -- node[edge label, left] {$6$} (a);
				
				\draw[->] (a) -- node[edge label, sloped, above, pos=0.65] {$16$} (d);
				
				\draw[->] (d) -- node[edge label, right] {$2$} (b);
				\draw[->] (d) -- node[edge label, below] {$2$} (c);
				
				\draw[->] (b) -- node[edge label, sloped, below, pos=0.65] {$4$} (c);				
			\end{tikzpicture}
			
			\caption{Graf większościowy dla drugiego profilu preferencji (po dodaniu ośmiu wyborców)}
			\label{fig:drugi_profil_graf}
			
		\end{figure}
		
		W nowym profilu uczestniczy \(32\) wyborców. Otrzymujemy:
		\[
		m_c = -4
		\qquad \text{oraz} \qquad
		g_R(c,a)=6,
		\]
		więc — ponownie na mocy Lematu~\ref{lem:condorcet_obecnosc} — metoda \(M\) nie może wybrać kandydata \(a\).
		
		Ponadto:
		\[
		m_b=-2
		\qquad \text{oraz} \qquad
		g_R(b,c)=4,
		\]
		więc metoda \(M\) nie może wybrać kandydata \(c\).
		
		Jednak pozytywne kryterium obecności wymaga, aby metoda \(M\) wybrała \(a\) lub \(c\). 
		
		Zatem przed zmianą profilu zbiór zwycięzców jest zawarty w \(\{a,c\}\).
		Po zmianie profilu żaden z tych kandydatów nie może być wybrany, co przeczy pozytywnemu kryterium obecności.
	\end{proof}
	
	\begin{proof}
		Załóżmy, że zbiór wyborców \(W\) zawiera co najmniej \(34\) wyborców oraz że metoda \(M\) spełnia negatywne kryterium obecności.
		
		Rozważmy czterech kandydatów:
		\[
		a,b,c,d.
		\]
		
		Niech \(u\) oznacza profil preferencji przedstawiony w Tabeli~\ref{tab:profil_v_typy}. W szczególności każdy kandydat \(x \in A \setminus \{a,b,c,d\}\) znajduje się w preferencjach wszystkich wyborców poniżej kandydatów \(a,b,c,d\), zatem nie wpływa on na relacje większościowe pomiędzy tymi czterema kandydatami.
		
		\begin{table}[h]
			\centering
			\[
			\begin{array}{|c|c|c|c|c|c|}
				\hline
				\alpha & \beta & \gamma & \delta & \varepsilon & \zeta \\
				\hline
				4 & 5 & 1 & 6 & 4 & 6 \\
				\hline
				b & a & d & a & d & b \\
				c & c & a & d & b & d \\
				a & d & b & b & c & c \\
				dx & bx & cx & cx & ax & ax \\
				\hline
			\end{array}
			\]
			\caption{Profil preferencji $R^W$}
			\label{tab:profil_v_typy}
		\end{table}
		
		Odpowiadający profilowi \(u\) graf większościowy (dla kandydatów \(a,b,c,d\)) przedstawiono na Rysunku~\ref{fig:graf_u}.
		
		\begin{figure}[h]
			\centering
			
			\begin{tikzpicture}[
				>=stealth,
				thick,
				vertex/.style={circle, draw, minimum size=1.2cm},
				edge label/.style={draw=none, fill=none, inner sep=1pt}
				]
				
				% Wierzchołki
				\node[vertex] (a) at (0,6) {$a$};
				\node[vertex] (b) at (6,6) {$b$};
				\node[vertex] (c) at (0,0) {$c$};
				\node[vertex] (d) at (6,0) {$d$};
				
				% relacje większościowe
				\draw[->] (b) -- node[edge label, above] {$2$} (a);
				\draw[->] (c) -- node[edge label, left] {$2$} (a);
				
				\draw[->] (a) -- node[edge label, sloped, above, pos=0.65] {$4$} (d);
				
				\draw[->] (d) -- node[edge label, right] {$6$} (b);
				\draw[->] (d) -- node[edge label, below] {$8$} (c);
				
				\draw[->] (b) -- node[edge label, sloped, below, pos=0.65] {$16$} (c);
				
			\end{tikzpicture}
			
			\caption{Graf większościowy dla profilu \(R^W\)}
			\label{fig:graf_u}
			
		\end{figure}
		
		W profilu \(R^W\) uczestniczy \(26\) wyborców. Rozważmy zmiany profilu poprzez usuwanie wybranych wyborców:
		
		- po usunięciu trzech wyborców typu \(\alpha\) kandydat \(a\) staje się zwycięzcą Condorceta,
		a więc na mocy negatywnego kryterium obecności metoda \(M\) nie może wykluczyć kandydata \(a\) w profilu \(R^W\);
		
		- po usunięciu pięciu wyborców typu \(\beta\) kandydat \(d\) staje się zwycięzcą Condorceta, natomiast kandydat \(b\), umieszczany przez tych wyborców na ostatnim miejscu, nie może zostać wybrany w profilu \(R^W\).
		
		W konsekwencji zbiór zwycięzców w profilu \(R^W\) zawiera się w:
		\[
		\{a,c\}.
		\]
		
		Dodajmy teraz ośmiu nowych wyborców, otrzymując profil \(R^V\). Każdy z nich preferuje:
		\[
		c \succ a \succ b \sim d  \succ x \quad \text{dla } x \in A \setminus \{a,b,c,d\}.
		\]
		
		\begin{table}[h]
			\centering
			\[
			\begin{array}{|c|c|c|c|c|c|c|}
				\hline
				\alpha & \beta & \gamma & \delta & \varepsilon & \zeta & \eta \\
				\hline
				4 & 5 & 1 & 6 & 4 & 6 & 8 \\
				\hline
				b & a & d & a & d & b & c \\
				c & c & a & d & b & d & a \\
				a & d & b & b & c & c & b\sim d \\
				d & b & c & c & a & a & x \\
				\hline
			\end{array}
			\]
			\caption{Profil preferencji \(R^V\)}
			\label{tab:profil_v_typy}
		\end{table}
		
		Odpowiadający profilowi \(R^V\) graf większościowy przedstawiono na Rysunku~\ref{fig:graf_v}.
		
		\begin{figure}[h]
			\centering
			
			\begin{tikzpicture}[
				>=stealth,
				thick,
				vertex/.style={circle, draw, minimum size=1.2cm},
				edge label/.style={draw=none, fill=none, inner sep=1pt}
				]
				
				% Wierzchołki
				\node[vertex] (a) at (0,6) {$a$};
				\node[vertex] (b) at (6,6) {$b$};
				\node[vertex] (c) at (0,0) {$c$};
				\node[vertex] (d) at (6,0) {$d$};
				
				% relacje większościowe
				\draw[->] (a) -- node[edge label, above] {$6$} (b);
				\draw[->] (c) -- node[edge label, left] {$10$} (a);
				
				\draw[->] (a) -- node[edge label, sloped, above, pos=0.65] {$12$} (d);
				
				\draw[->] (d) -- node[edge label, right] {$6$} (b);
				
				\draw[->] (b) -- node[edge label, sloped, below, pos=0.65] {$8$} (c);
				
			\end{tikzpicture}
			
			\caption{Graf większościowy dla profilu \(R^V\)}
			\label{fig:graf_v}
			
		\end{figure}
		
		W nowym profilu uczestniczy \(34\) wyborców. Ponownie rozważamy operacje usuwania wyborców:
		
		- po usunięciu czterech wyborców typu \(\varepsilon\) oraz pięciu spośród sześciu wyborców typu \(\zeta\), kandydat \(c\) staje się zwycięzcą Condorceta, zatem metoda \(M\) nie może wybrać kandydata \(a\) w profilu \(R^V\);
		
		- po usunięciu jednego wyborcy typu \(\gamma\) oraz sześciu wyborców typu \(\delta\), kandydat \(b\) staje się zwycięzcą Condorceta, zatem metoda \(M\) nie może wybrać kandydata \(c\) w profilu \(R^V\).
		
		Otrzymujemy więc, że żaden z kandydatów \(a\) i \(c\) nie może zostać wybrany w profilu \(v\), co przeczy negatywnemu kryterium obecności, które wymaga, aby co najmniej jeden z nich należał do zbioru zwycięzców.
		
		To kończy dowód.
	\end{proof}	
	
	\begin{theorem}[Moulin \cite{moulin1988}]
		Jeżeli liczba kandydatów spełnia \(\#K \geq 4\) oraz liczba wyborców wynosi conajlminej 37 (34 w przypadku negatywnego kryterium obecności), to każda metoda wyborcza spełniająca kryterium Condorceta jest podatna na paradoks nieobecności.
	\end{theorem}
	
	\begin{proof}
		Korzystając z poprzednich twierdzeń oraz z faktu że brak spełnienia pozytywnego(negatywnego) kryterium obecności implikuje brak spełnienia kryterium obecności dostajemy teze.
	\end{proof}
	
	
	
	
	
	
	
	
	
	

	\section{Metody nie spełniające kryterium Condorceta}
	
	Twierdzenie to pokazuje fundamentalne napięcie pomiędzy naturalnymi własnościami metod wyborczych: zgodność z kryterium Condorceta implikuje możliwość wystąpienia paradoksu nieobecności.
	
	\medskip
	
	Warto zauważyć, że implikacja odwrotna nie zachodzi. Istnieją metody wyborcze, które nie spełniają kryterium Condorceta, a mimo to są podatne na paradoks nieobecności (np. metoda Hare'a). Z drugiej strony istnieją metody odporne na ten paradoks, takie jak dyktatura czy metoda punktów Bordy.
	
	\begin{proof}[Dyktatura jest metodą odporną na nieobecność]
		
		Niech \(M\) będzie metodą dyktatorską z dyktatorem \(d \in W\).
		
		Z definicji metody dyktatorskiej wynik wyborów jest zawsze równy najbardziej preferowanemu kandydatowi dyktatora. Oznacza to, że dla każdego profilu preferencji \(R\)
		\[
		M(R)=\max_{\overset{d}{\succ}} K.
		\]
		
		Rozważmy dowolnego wyborcę \(v \in W\).
		
		\begin{itemize}
			\item Jeżeli \(v=d\), to po usunięciu dyktatora wynik wyborów albo pozostaje taki sam, albo zmienia się na mniej preferowany z punktu widzenia dyktatora. W szczególności
			\[
			M(R) \overset{d}{\succsim} M(R-d).
			\]
			
			\item Jeżeli \(v\neq d\), to usunięcie wyborcy \(v\) nie wpływa na wynik wyborów, ponieważ wynik zależy wyłącznie od preferencji dyktatora. Zatem
			\[
			M(R) \overset{v}{=} M(R-v).
			\]
		\end{itemize}
		
		W obu przypadkach wyborca nie uzyskuje lepszego wyniku poprzez rezygnację z udziału w głosowaniu.
		
		W konsekwencji metoda dyktatorska jest odporna na paradoks nieobecności.
	\end{proof}
	
	\begin{proof}[Metoda punktów Bordy jest odporna na paradoks nieobecności]
		Niech \(R\) będzie profilem preferencji oraz \(w \in W\) dowolnym wyborcą.
		
		Bez straty ogólności załóżmy, że preferencje wyborcy \(w\) mają postać
		\[
		k_1 \overset{w}{\succ} k_2 \overset{w}{\succ} \dots \overset{w}{\succ} k_m.
		\]
		
		W metodzie Bordy kandydat znajdujący się na \(i\)-tej pozycji otrzymuje \(m-i\) punktów.
		
		Niech
		\[
		s_i
		\]
		oznacza liczbę punktów uzyskaną przez kandydata \(k_i\) w profilu \(R\).
		
		Po usunięciu wyborcy \(w\) kandydat \(k_i\) traci dokładnie \(m-i\) punktów, a więc jego nowy wynik wynosi
		\[
		s_i-(m-i).
		\]
		
		Zauważmy, że kandydaci wyżej oceniani przez wyborcę \(w\) tracą więcej punktów niż kandydaci oceniani niżej.
		
		W szczególności, jeżeli
		\[
		i<j,
		\]
		to
		\[
		(m-i) > (m-j).
		\]
		
		Oznacza to, że usunięcie wyborcy \(w\) działa relatywnie bardziej na korzyść kandydatów mniej preferowanych przez \(w\).
		
		W konsekwencji wynik wyborów po usunięciu wyborcy \(w\) nie może być przez niego ściśle preferowany względem wyniku przy jego udziale, tj.
		\[
		M(R-w)\overset{w}{\leq} M(R).
		\]
		
		Zatem metoda punktów Bordy spełnia kryterium obecności.
	\end{proof}
	
	\begin{example}[Metoda Hare'a a paradoks nieobecności]
		Rozważmy metodę Hare'a (IRV) oraz profil preferencji przedstawiony w tabeli:
		\[
		\begin{array}{|c|c|c|c|}
			\hline
			10 & 5 & 4 & 6\\
			\hline
			a & b & b & c\\
			b & c & c & d\\
			c & a & d & a\\
			d & d & a & b\\
			\hline
		\end{array}
		\]
		
		Dla powyższego profilu metoda Hare'a wskazuje jako zwycięzcę kandydata \(A\).
		
		Rozważmy teraz sytuację, w której grupa \(4\) wyborców o preferencjach
		\[
		b \succ c \succ d \succ a
		\]
		rezygnuje z udziału w głosowaniu. Otrzymujemy wówczas zmodyfikowany profil:
		\[
		\begin{array}{|c|c|c|}
			\hline
			10 & 5 & 6\\
			\hline
			a & b & c\\
			b & c & d\\
			c & a & a\\
			d & d & b\\
			\hline
		\end{array}
		\]
		
		W tym przypadku zwycięzcą zostaje kandydat \(C\).
		
		Ponieważ dla wyborców o preferencjach \(b \succ c \succ d \succ a\) zachodzi
		\[
		c \succ a,
		\]
		oznacza to, że wynik po ich rezygnacji z udziału w głosowaniu jest przez nich ściśle preferowany względem wyniku uzyskanego przy ich uczestnictwie.
		
		W konsekwencji metoda Hare'a jest podatna na paradoks nieobecności, mimo że nie spełnia kryterium Condorceta.
	\end{example}
	
	Przykład ten pokazuje, że brak spełnienia kryterium Condorceta nie implikuje odporności na paradoks nieobecności.
	
	\section{Inne twierdzenia}
	
	[TODO - cała sekcja, póki co ta część jest tylko przetłumaczona]
	
	\begin{lemma}
		\label{lemma:sat_1}
		Załóżmy, że metoda wyborcza \(M\) spełnia kryterium Condorceta oraz kryterium obecności. Niech \(R\) będzie profilem preferencji, a \(B \subseteq K\) zbiorem tzw. \emph{złych} alternatyw, przy czym każdy wyborca klasyfikuje każdą alternatywę \(x \in B\) niżej niż każdą alternatywę \(y \in K \setminus B\). Wówczas
		\[
		M(R) \notin B.
		\]
	\end{lemma}
	
	\begin{proof}
		Dowód przeprowadzimy indukcyjnie względem liczby wyborców.
		
		\textbf{1 krok indukcyjny} Załóżmy, że elektorat składa się z jednego wyborcy, tj. \(R=\{(w,\overset{w}{\succ})\}\). Wówczas kandydatem Condorceta jest alternatywa znajdująca się na pierwszym miejscu w preferencjach wyborcy \(w\). Z kryterium Condorceta wynika więc, że metoda \(M\) wskazuje tę właśnie alternatywę jako zwycięzcę. Ponieważ każda alternatywa z \(B\) jest klasyfikowana niżej niż każda alternatywa z \(K\setminus B\), zwycięzca nie należy do \(B\).
		
		\textbf{2 krok indukcyjny} Załóżmy, że teza zachodzi dla każdego profilu złożonego z \(n-1\) wyborców, gdzie \(n \geq 2\), i rozważmy profil \(R\) złożony z \(n\) wyborców. Niech \(v \in V\) będzie dowolnym wyborcą.
		
		Z warunku uczestnictwa otrzymujemy
		\[
		M(R) \overset{v}{\succsim} M(R-v).
		\]
		
		Przypuśćmy nie wprost, że \(M(R)\in B\). Ponieważ każdy kandydat spoza zbioru \(B\) jest przez wyborcę \(v\) oceniany wyżej niż każdy kandydat należący do \(B\), z powyższej nierówności wynika, że również \(M(R-v)\in B\).
		
		Profil \(R-v\) ma jednak \(n-1\) wyborców, zatem z założenia indukcyjnego otrzymujemy
		\[
		M(R-v)\notin B,
		\]
		co prowadzi do sprzeczności.
		
		Wobec tego
		\[
		M(R)\notin B.
		\]
		
		Indukcja kończy dowód.
	\end{proof}
	
	\begin{theorem}
		Nie istnieje metoda wyborcza spełniająca kryterium Maximina, która jednocześnie spełnia kryterium obecności dla
		\[
		\#K \geq 4
		\quad \text{oraz} \quad
		\#W \geq 7.
		\]
	\end{theorem}
	
	\begin{proof}
		Niech \(M\) będzie metodą spełniającą kryterium Maximina oraz warunek uczestnictwa.
		
		Rozważmy profil sześciu wyborców \(R\) przedstawiony w Tabeli~\ref{fig:profil_R}.
		
		\begin{table}[h]
			\centering
			\[
			\begin{array}{|c|c|c|c|}
				\hline
				1 & 2 & 2 & 1\\
				\hline
				a & b & c & d\\
				b & d & a & c\\
				d & c & b & a\\
				c & a & d & b\\
				\hline
			\end{array}
			\]
			\caption{Profil preferencji \(R\)}
			\label{fig:profil_R}
		\end{table}
		
		Przypuśćmy, że \(M(R) \in \{a,b\}\). Dodanie głosu typu \(abcd\) prowadzi do ważonego turnieju, w którym alternatywa \(c\) jest jedynym zwycięzcą Maximina, zatem metoda \(M\) musi wskazać \(c\) jako zwycięzce. Prowadzi to jednak do sprzeczności z kryterium obecności, ponieważ wyborca, który dodał swój głos, uzyskałby lepszy wynik, gdyby wstrzymał się od głosowania.
		
		\begin{table}[h]
			\centering
			\[
			\begin{array}{|c|c|c|c|c|}
				\hline
				1 & 2 & 2 & 1 & 1\\
				\hline
				a & b & c & d & a\\
				b & d & a & c & b\\
				d & c & b & a & c\\
				c & a & d & b & d\\
				\hline
			\end{array}
			\]
			\caption{Profil otrzymany po dodaniu głosu typu \(abcd\)}
			\label{tab:profil_abcd}
		\end{table}
	
		\begin{figure}[h]
			\centering
			
			\begin{tikzpicture}[
				>=stealth,
				thick,
				vertex/.style={circle, draw, minimum size=1.2cm},
				edge label/.style={draw=none, fill=none, inner sep=1pt}
				]
				
				% Wierzchołki
				\node[vertex] (a) at (0,6) {$a$};
				\node[vertex] (b) at (6,6) {$b$};
				\node[vertex] (c) at (0,0) {$c$};
				\node[vertex] (d) at (6,0) {$d$};
				
				% relacje większościowe
				\draw[->] (a) -- node[edge label, above] {$3$} (b);
				\draw[->] (c) -- node[edge label, left] {$3$} (a);
				
				\draw[->] (a) -- node[edge label, sloped, above, pos=0.65] {$1$} (d);
				
				\draw[->] (b) -- node[edge label, right] {$5$} (d);
				\draw[->] (d) -- node[edge label, below] {$1$} (c);
				
				\draw[->] (b) -- node[edge label, sloped, below, pos=0.65] {$1$} (c);
				
			\end{tikzpicture}
			
			\caption{Graf otrzymany po dodaniu głosu typu \(abcd\)}
			\label{fig:graf_1_1}
			
		\end{figure}
		
		Analogicznie, jeśli \(M(R) \in \{c,d\}\), to dodanie głosu typu \(dcba\) prowadzi do ważonego turnieju, w którym \(b\) jest jedynym zwycięzcą Maximina, co ponownie przeczy kryterium obecności.
		
		\begin{table}[h]
			\centering
			\[
			\begin{array}{|c|c|c|c|c|}
				\hline
				1 & 2 & 2 & 1 & 1\\
				\hline
				a & b & c & d & d\\
				b & d & a & c & c\\
				d & c & b & a & b\\
				c & a & d & b & a\\
				\hline
			\end{array}
			\]
			\caption{Profil otrzymany po dodaniu głosu typu \(dcba\)}
			\label{tab:profil_dcba}
		\end{table}
	
		\begin{figure}[h]
			\centering
			
			\begin{tikzpicture}[
				>=stealth,
				thick,
				vertex/.style={circle, draw, minimum size=1.2cm},
				edge label/.style={draw=none, fill=none, inner sep=1pt}
				]
				
				% Wierzchołki
				\node[vertex] (a) at (0,6) {$a$};
				\node[vertex] (b) at (6,6) {$b$};
				\node[vertex] (c) at (0,0) {$c$};
				\node[vertex] (d) at (6,0) {$d$};
				
				% relacje większościowe
				\draw[->] (b) -- node[edge label, above] {$2$} (a);
				\draw[->] (c) -- node[edge label, left] {$2$} (a);
				
				\draw[->] (a) -- node[edge label, sloped, above, pos=0.65] {$4$} (d);
				
				\draw[->] (d) -- node[edge label, right] {$6$} (b);
				\draw[->] (d) -- node[edge label, below] {$8$} (c);
				
				\draw[->] (b) -- node[edge label, sloped, below, pos=0.65] {$16$} (c);
				
			\end{tikzpicture}
			
			\caption{Graf otrzymany po dodaniu głosu typu \(dcba\)}
			\label{fig:graf_1_2}
			
		\end{figure}
		
		Symetria tego argumentu wynika z automorfizmu profilu \(R\), polegającego na przekształceniu etykiet alternatyw zgodnie z odwzorowaniem \(abcd \mapsto dcba\).
		
		Jeżeli \(\#K > 4\), do każdego profilu wyborców w \(R\) oraz do wszystkich dodatkowych wyborców użytych w dowodzie dodajemy nowe „złe” alternatywy \(x_1, x_2, \dots, x_m\) umieszczone na dole każdego rankingu. Z Lematu~~\ref{lemma:sat_1} wynika, że \(M\) wybiera w każdym kroku dowodu spośród zbioru \(\{a,b,c,d\}\), więc dowód przebiega identycznie jak wcześniej.
	\end{proof}
	
	Wynik istnienia dla \(n \leq 6\) został uzyskany metodami opisanymi w~\cite{brandt2017sat}.
	
	\begin{theorem}
		Nie istnieje metoda wyborcza spełniająca kryterium Kemeny'ego, która jednocześnie spełnia kryterium obecności dla
		\[
		\#K \geq 4
		\quad \text{oraz} \quad
		\#W \geq 4.
		\]
	\end{theorem}
	
	\begin{proof}
		Niech \(M\) będzie metodą wyborczą spełniającą kryterium Kemeny'ego oraz kryterium obecności.
		
		Rozważmy profil czterech wyborców \(R\) przedstawiony w Tabeli~\ref{tab:kemeny_1}.
		
		\begin{table}[h]
			\centering
			\[
			\begin{array}{|c|c|c|c|}
				\hline
				1 & 1 & 1 & 1\\
				\hline
				a & d & c & b\\
				b & a & d & c\\
				c & b & a & d\\
				d & c & b & a\\
				\hline
			\end{array}
			\]
			\caption{Profil preferencji \(R\)}
			\label{tab:kemeny_1}
		\end{table}
	
		\begin{table}[h]
			\centering
			\[
			\begin{array}{c|cccc}
				\Rsh & a & b & c & d\\
				\hline
				a & - & 0 & 0 & 0\\
				b & 0 & - & 0 & 0\\
				c & 0 & 0 & - & 0\\
				d & 0 & 0 & 0 & -\\
			\end{array}
			\]
			\caption{Macierz marginesów większości dla profilu}
			\label{tab:marginesy_kemeny}
		\end{table}
		
		Przypuśćmy, że \(M(R)=c\). Wówczas usunięcie wyborcy z preferencji \(cbad\) z profilu \(R\) prowadzi do ważonego turnieju, w którym ranking \(adcb\) jest jedynym rankingiem Kemeny'ego (o wartości 12), a zatem kandydat \(a\) jest jedynym zwycięzcą Kemeny'ego. Otrzymujemy sprzeczność z kryterium obecności.
		
		Analogicznie można wykluczyć pozostałe trzy możliwe wyniki dla profilu \(R\), rozważając sytuację, w której jeden z wyborców się wstrzymuje od głosowania. W każdym przypadku otrzymujemy jednoznacznego zwycięzcę Kemeny'ego oraz sprzeczność z kryterium obecności.
		
		Argumenty te są identyczne, ponieważ profil \(R\) jest całkowicie symetryczny w tym sensie, że dla każdej pary kandydatów \(x\) i \(y\) istnieje automorfizm profilu \(R\), który odwzorowuje \(x\) na \(y\).
		
		Podobnie jak w Twierdzeniu~1, jeżeli \(\#K > 4\), do profilu \(R\) oraz do dodatkowych wyborców dodajemy nowych kandydatów \(x_1, x_2, \dots, x_{\#K-4}\), umieszczone na dole każdego rankingu. Z Lematu~1 wynika, że \(M\) wybiera w każdym kroku spośród zbioru \(\{a,b,c,d\}\), co kończy dowód.
	\end{proof}
	
	\begin{theorem}
		Nie istnieje metoda wyborcza spełniająca kryterium Condorceta, która jednocześnie spełnia kryterium obecności dla
		\[
		|K| \geq 4
		\quad \text{oraz} \quad
		|W| \geq 12.
		\]
	\end{theorem}
	
	\begin{proof}
		Rozważmy najpierw przypadek \(m=4\). Dowód ma strukturę przedstawioną na Rysunku~3. Niech \(R\) będzie profilem preferencji pokazanym na tym rysunku.
		
		Ponieważ profil \(R\) pozostaje niezmieniony po przestawieniu etykiet alternatyw zgodnie z odwzorowaniem \(abcd \mapsto dcba\) oraz po permutacji wyborców, możemy bez straty ogólności założyć, że \(f(R) \in \{a,b\}\). (Jawny dowód dla przypadku \(f(R) \in \{c,d\}\) został zaznaczony na Rysunku~3.)
		
		Z warunku uczestnictwa wynika, że z \(f(R) \in \{a,b\}\) mamy również \(f(R_\alpha := R + 2\cdot abcd) \in \{a,b\}\), ponieważ wyborcy o preferencjach \(abcd\) nie mogą pogorszyć swojego wyniku przez dołączenie do elektoratu.
		
		Jeżeli \(f(R_\alpha)=a\), to ponownie z warunku uczestnictwa wynika, że usunięcie dwóch wyborców o preferencjach \(bdca\) nie zmienia zwycięzcy, tj. \(f(R_\alpha - 2\cdot bdca)=a\). Dodanie profilu \(acdb\) również nie zmienia wyniku, więc
		\[
		f(R_\alpha - 2\cdot bdca + acdb)=a.
		\]
		Jest to jednak sprzeczne z faktem, że profil \(R_\alpha - 2\cdot bdca + acdb\) posiada zwycięzcę Condorceta, którym jest \(c\).
		
		Zatem musi zachodzić \(f(R_\alpha)=b\), co implikuje \(f(R_\alpha - dcab)=b\), a więc
		\[
		f(R_\beta := R_\alpha - dcab - 2\cdot cabd) \in \{b,d\}.
		\]
		
		Rozpatrujemy ponownie przypadki:
		
		Jeżeli \(f(R_\beta)=b\), to możemy dodać dwóch wyborców o preferencjach \(badc\), otrzymując profil z jednoznacznym zwycięzcą Condorceta \(a\), co przeczy warunkowi uczestnictwa.
		
		Jeżeli natomiast \(f(R_\beta)=d\), to otrzymujemy \(f(R_\beta - abcd)=d\), a następnie z warunku uczestnictwa
		\[
		f(R_\beta - abcd + 3\cdot dcba)=d,
		\]
		co stoi w sprzeczności z istnieniem zwycięzcy Condorceta \(b\). Otrzymujemy więc sprzeczność.
		
		Dla \(m>4\) dodajemy „złe” alternatywy \(x_1, x_2, \dots, x_{m-4}\) na dole każdego rankingu w profilu \(R\) oraz u wszystkich dodatkowych wyborców. Z Lematu~1 wynika, że \(f\) wybiera na każdym etapie spośród zbioru \(\{a,b,c,d\}\), co pozwala przeprowadzić dowód identycznie jak w przypadku \(m=4\).
	\end{proof}
	
	\begin{theorem}[Twierdzenie 4]
		Istnieje rozszerzenie Condorceta \(f\), które spełnia warunek uczestnictwa dla \(m=4\) oraz \(n \leq 11\). Ponadto \(f\) jest metodą parową, jest efektywne w sensie Pareto oraz stanowi uszczegółowienie cyklu górnego (top cycle).
	\end{theorem}
	
	Metoda \(f\) jest zadana w postaci tablicy wyników (look-up table), która została wyznaczona przy użyciu solvera SAT. Tablica ta zawiera wszystkie \(1{,}204{,}215\) ważonych turniejów indukowalnych przez co najwyżej 11 wyborców i przyporządkowuje każdemu z nich alternatywę zwycięską (zob. Rysunek~4 dla fragmentu tej tablicy).
	
	\begin{lemma}
		Niech \(F\) będzie wielowartościowym rozszerzeniem Condorceta. Niech \(R\) będzie profilem preferencji oraz \(B \subset A\) zbiorem tzw. „złych” alternatyw, przy czym każdy wyborca klasyfikuje każdą alternatywę \(x \in B\) niżej niż każdą alternatywę \(y \in A \setminus B\).
		
		Jeżeli \(F\) spełnia optymistyczny warunek uczestnictwa, to
		\[
		F(R) \not\subseteq B,
		\]
		tj. zbiór wyników zawiera przynajmniej jedną nie-złą alternatywę.
		
		Jeżeli \(F\) spełnia pesymistyczny warunek uczestnictwa, to
		\[
		F(R) \cap B = \emptyset,
		\]
		tj. żadna zła alternatywa nie należy do zbioru wyników.
	\end{lemma}
	
	\begin{proof}
		\textbf{Optymistyczny warunek uczestnictwa.}
		Dowód przeprowadzamy przez indukcję względem liczby wyborców \(|N|\) w profilu \(R\).
		
		Jeżeli \(R\) składa się z jednego wyborcy \(i\), to ponieważ \(F\) jest rozszerzeniem Condorceta, mamy \(F(R)=\{x\}\), gdzie \(x\) jest najlepszą alternatywą wyborcy \(i\), a więc nie jest alternatywą złą.
		
		Jeżeli \(R\) zawiera co najmniej dwóch wyborców i \(i \in N\), to z optymistycznego warunku uczestnictwa wynika, że
		\[
		\max_{\prec_i} F(R) \succ_i \max_{\prec_i} F(R - i),
		\]
		ponieważ \(R = R - i + \prec_i\).
		
		Z założenia indukcyjnego zbiór \(F(R - i)\) zawiera alternatywę nie-złą \(x\), zatem również \(\max_{\prec_i} F(R - i)\) nie jest alternatywą złą. Wynika stąd, że \(\max_{\prec_i} F(R)\) również nie jest zła, a więc \(F(R)\) zawiera alternatywę nie-złą.
		
		\textbf{Pesymistyczny warunek uczestnictwa.}
		Dowód prowadzimy przez indukcję względem \(|N|\).
		
		Jeżeli \(R\) składa się z jednego wyborcy \(i\), to \(F(R)=\{x\}\), gdzie \(x\) jest najlepszą alternatywą wyborcy \(i\), a więc nie jest alternatywą złą.
		
		Jeżeli \(R\) zawiera co najmniej dwóch wyborców i \(i \in N\), to z pesymistycznego warunku uczestnictwa mamy
		\[
		\min_{\prec_i} F(R) \succ_i \min_{\prec_i} F(R - i).
		\]
		
		Z założenia indukcyjnego zbiór \(F(R - i)\) nie zawiera alternatyw złych, zatem \(\min_{\prec_i} F(R - i)\) również nie jest alternatywą złą. Wynika stąd, że \(\min_{\prec_i} F(R)\) nie jest zła, a więc — z definicji minimum — zbiór \(F(R)\) nie zawiera żadnej alternatywy złej.
	\end{proof}
	
	\begin{theorem}[Twierdzenie 5]
		Istnieje wielowartościowe rozszerzenie Condorceta \(F\), które spełnia optymistyczny warunek uczestnictwa dla \(m=4\) oraz \(n \leq 16\), a ponadto jest efektywne w sensie Pareto oraz stanowi uszczegółowienie cyklu górnego (top cycle).
	\end{theorem}
	
	\begin{theorem}[Twierdzenie 6]
		Nie istnieje wielowartościowe rozszerzenie Condorceta, które spełnia optymistyczny warunek uczestnictwa dla \(m \geq 4\) oraz \(n \geq 17\).
	\end{theorem}
	
	\begin{proof}
		Niech \(F\) będzie taką funkcją i rozważmy profil 10 wyborców \(R\) przedstawiony na Rysunku~5.
		
		Przypuśćmy, że \(a \in F(R)\) lub \(b \in F(R)\). (Przypadek \(c \in F(R)\) lub \(d \in F(R)\) jest symetryczny.) Niech \(R_\alpha := R + 2\cdot abcd\). Z optymistycznego warunku uczestnictwa wynika, że również \(a \in F(R_\alpha)\) lub \(b \in F(R_\alpha)\).
		
		Jeżeli \(a \in F(R_\alpha)\), to z warunku uczestnictwa otrzymujemy także
		\[
		a \in F(R_\alpha + 3\cdot acbd),
		\]
		jednak profil ten posiada zwycięzcę Condorceta \(c\), co prowadzi do sprzeczności.
		
		Jeżeli natomiast \(b \in F(R_\alpha)\), to również
		\[
		b \in F(R_\alpha + 5\cdot bacd),
		\]
		lecz ten profil ma zwycięzcę Condorceta \(a\), co również prowadzi do sprzeczności.
		
		Dowód dla więcej niż 4 alternatyw wynika bezpośrednio z Lematu~2.
	\end{proof}
	
	\begin{theorem}[Twierdzenie 7]
		Nie istnieje wielowartościowe rozszerzenie Condorceta, które spełnia pesymistyczny warunek uczestnictwa dla \(m \geq 4\) oraz \(n \geq 14\). Z drugiej strony, dla \(m=4\) oraz \(n \leq 13\) istnieje taka wielowartościowa reguła wyborcza.
	\end{theorem}
	
	\begin{proof}
		Dowód ma podobną strukturę jak dowód Twierdzenia~3 i został przedstawiony w formie graficznej na Rysunku~6.
		
		Początkowy profil w tym dowodzie ma postać
		\[
		R^* := R + 2\cdot abcd + 2\cdot dcba,
		\]
		gdzie \(R\) jest profilem z Rysunku~3. W dalszej części zastępujemy kroki, w których dodawani są wyborcy, analogicznymi krokami, w których wyborcy są usuwani, i w każdym kroku stosujemy pesymistyczny warunek uczestnictwa, aby otrzymać sprzeczność.
		
		Ponieważ profil \(R^*\) pozostaje niezmieniony po relabelingu alternatyw zgodnie z \(abcd \mapsto dcba\) oraz permutacji wyborców, możemy bez straty ogólności założyć, że \(a \in F(R^*)\) lub \(b \in F(R^*)\).
		
		Niech \(R_\alpha := R^* - 2\cdot dcba\). Z pesymistycznego warunku uczestnictwa wynika, że również \(a \in F(R_\alpha)\) lub \(b \in F(R_\alpha)\), ponieważ wyborcy o preferencjach \(abcd\) nie powinni być lepiej sytuowani po opuszczeniu elektoratu.
		
		Jeżeli \(a \in F(R_\alpha)\), to z warunku uczestnictwa mamy także
		\[
		a \in F(R_\alpha - 3\cdot bdca),
		\]
		jednak profil \(R_\alpha - 3\cdot bdca\) posiada zwycięzcę Condorceta \(c\), co prowadzi do sprzeczności.
		
		Zatem musi zachodzić \(b \in F(R_\alpha)\), co implikuje \(b \in F(R_\alpha - dcab)\). Niech
		\[
		R_\beta := R_\alpha - dcab - 3\cdot cabd.
		\]
		Wówczas z warunku uczestnictwa wynika, że \(F(R_\beta)\) zawiera \(b\) lub \(d\).
		
		Rozpatrujemy przypadki. Jeżeli \(b \in F(R_\beta)\), to usunięcie wyborcy \(dcab\) prowadzi do profilu z jednoznacznym zwycięzcą Condorceta \(a\), co przeczy warunkowi uczestnictwa.
		
		Jeżeli natomiast \(d \in F(R_\beta)\), to mamy \(d \in F(R_\beta - 2\cdot abcd)\), a następnie z warunku uczestnictwa
		\[
		F(R_\beta - 2\cdot abcd - abdc)
		\]
		zawiera \(c\) lub \(d\), co stoi w sprzeczności z istnieniem zwycięzcy Condorceta \(b\). Otrzymujemy więc sprzeczność.
	\end{proof}
	
	\begin{theorem}[Twierdzenie 8]
		Nie istnieje wielowartościowe rozszerzenie Condorceta, które jednocześnie spełnia optymistyczny oraz pesymistyczny warunek uczestnictwa dla \(m \geq 4\) oraz \(n \geq 12\). Z drugiej strony, dla \(m=4\) oraz \(n \leq 11\) istnieje taka wielowartościowa reguła wyborcza (która dodatkowo jest efektywna w sensie Pareto oraz stanowi uszczegółowienie cyklu górnego).
	\end{theorem}
	
	\begin{proof}
		Wykorzystujemy dowód Twierdzenia~3, przy czym stosujemy optymistyczny warunek uczestnictwa dla krawędzi oznaczonych dodaniem wyborcy \((+)\), natomiast pesymistyczny warunek uczestnictwa dla krawędzi oznaczonych usunięciem wyborcy \((-)\).
		
		Z drugiej strony, jednoznaczna (rozstrzygająca) reguła wyborcza z Twierdzenia~4 spełnia zarówno optymistyczny, jak i pesymistyczny warunek uczestnictwa.
	\end{proof}
	
	
	
	
	
	
	
	
	
	
	
	
	
	
	
	
	
	
	\chapter{Silny paradoks nieobecności}
	
	Klasyczny paradoks nieobecności polega na tym, że pewna grupa wyborców może uzyskać bardziej preferowany wynik poprzez rezygnację z udziału w głosowaniu. W wielu przypadkach zjawisko to nie wydaje się jednak szczególnie istotne, jeżeli uzyskana w ten sposób poprawa dotyczy alternatywy, która nie jest dla tej grupy zdecydowanie preferowana.
	
	Prowadzi to do naturalnego wzmocnienia rozważanego zjawiska.
	
	\begin{definition}[silny paradoks nieobecności]
		Mówimy, że metoda wyborcza \(M\) jest podatna na silny paradoks nieobecności, jeżeli istnieje profil preferencji \(R \in \mathcal{R}^{V}\), grupa wyborców \(S \subseteq V\) oraz kandydat \(a \in K\), taki że:
		\begin{itemize}
			\item kandydat \(a\) jest najbardziej preferowaną alternatywą dla wszystkich wyborców z \(S\),
			\item przy pełnym udziale \(S\) w głosowaniu wynik \(M(R)\) jest przez nich mniej preferowany niż \(a\),
			\item natomiast po rezygnacji grupy \(S\) z udziału w głosowaniu zachodzi
			\[
			a \succ M(R - S).
			\]
		\end{itemize}
	\end{definition}
	
	Innymi słowy, w silnej wersji paradoksu nieobecności grupa wyborców może doprowadzić do wyboru swojej najbardziej preferowanej alternatywy poprzez wstrzymanie się od głosu, mimo że aktywny udział w głosowaniu prowadzi do wyniku dla niej gorszego.
	
	\medskip
	
	Zjawisko to można zilustrować na przykładzie procedury głosowania sekwencyjnego.
	
	\begin{example}[Silny paradoks nieobecności – procedura sekwencyjna]
		Rozważmy profil preferencji:
		\[
		\begin{array}{|c|c|c|c|}
			\hline
			2 & 3 & 2 & 2\\
			\hline
			A & B & C & C\\
			B & C & A & B\\
			C & A & B & A\\
			\hline
		\end{array}
		\]
		
		Przy ustalonej agendzie głosowań:
		\[
		A \text{ vs } B,\quad \text{następnie zwycięzca vs } C,
		\]
		wynikiem końcowym przy pełnej frekwencji jest kandydat \(B\).
		
		Rozważmy teraz sytuację, w której dwóch wyborców o preferencjach
		\[
		C \succ A \succ B
		\]
		rezygnuje z udziału w głosowaniu. Wówczas zmodyfikowany profil prowadzi do zwycięstwa kandydata \(C\), który jest dla tej grupy najbardziej preferowaną alternatywą.
		
		Oznacza to, że rezygnacja z udziału w głosowaniu prowadzi do wyniku bardziej korzystnego dla tych wyborców niż ich aktywny udział.
	\end{example}
	
	\medskip
	
	Silny paradoks nieobecności może występować nie tylko w procedurach zależnych od agendy, ale również w metodach bezpośrednich, takich jak metoda Coombsa czy metoda Nansona.
	
	W metodzie Coombsa oraz Nansona możliwe jest wystąpienie sytuacji, w której przy pełnej frekwencji wygrywa kandydat \(B\), natomiast po rezygnacji części wyborców zwycięzcą zostaje kandydat \(C\), który jest dla tej grupy bardziej preferowany.
	
	\medskip
	
	Zjawisko silnego paradoksu nieobecności prowadzi do naturalnego pytania o jego związek z kryteriami racjonalności metod wyborczych, w szczególności z kryterium Condorceta.
	
	\begin{theorem}[Moulin]
		Jeżeli \(\#K \geq 4\), to każda metoda wyborcza spełniająca kryterium Condorceta jest podatna na silny paradoks nieobecności.
	\end{theorem}
	
	Twierdzenie to pokazuje, że zgodność z kryterium Condorceta nie eliminuje możliwości wystąpienia silnego paradoksu nieobecności.
	
	\medskip
	
	W literaturze rozważane są również metody odporne na to zjawisko. Przykładowo metoda maximina nie wykazuje podatności na silny paradoks nieobecności.
	
	\medskip
	
	Istotnym pytaniem jest również istnienie warunków strukturalnych, które pozwalają kontrolować możliwość wystąpienia tego typu paradoksów. Jednym z podejść jest wprowadzenie ograniczeń typu \(q\)-majority.
	
	Niech \(q \in \left[\tfrac{1}{2},1\right]\). Mówimy, że kandydat \(a\) dominuje kandydata \(b\) w sensie \(q\)-większości, jeżeli
	\[
	\#\{w \in V : a \overset{w}{\succ} b\} \ge q \cdot \#V.
	\]
	
	Zbiór kandydatów niedominowanych definiujemy jako
	\[
	C(q,R) = \{a \in K \mid \nexists b \in K \text{ takie, że } b \text{ dominuje } a\}.
	\]
	
	\begin{definition}[\(q\)-warunek zgodności]
		Metoda wyborcza \(M\) spełnia warunek \(q\)-zgodności, jeżeli dla każdego profilu preferencji \(R\) zachodzi:
		\[
		C(q,R) \neq \emptyset \;\Rightarrow\; M(R) \in C(q,R).
		\]
	\end{definition}
	
	W tym kontekście można sformułować następujący wynik.
	
	\begin{theorem}[Holzman]
		Dla zbioru kandydatów \(K\) warunek \(q\)-zgodności oraz odporność na silny paradoks nieobecności są zgodne wtedy i tylko wtedy, gdy
		\[
		q \ge \frac{\#K - 1}{\#K}
		\quad \text{lub} \quad \#K \le 3.
		\]
	\end{theorem}
	
	Wynik ten pokazuje, że im większa liczba kandydatów, tym bardziej restrykcyjne warunki większościowe są konieczne, aby uniknąć występowania silnego paradoksu nieobecności.
	
	
	
	
	
	
	
	\chapter*{Podsumowanie}
	
	\newpage
	
	
	\begin{thebibliography}{99}
		
		\bibitem{brandt2017nsp}
		F.~Brandt, J.~Hofbauer, and M.~Strobel,
		\textit{Exploring the No-Show Paradox for Condorcet Extensions},
		in: M.~Diss and V.~Merlin (eds.),
		\textit{Evaluating Voting Systems with Probability Models:
			Essays by and in Honor of William Gehrlein and Dominique Lepelley},
		Springer, Studies in Social Choice and Welfare, 2020.
		
		\bibitem{brandt2017sat}
		F.~Brandt, C.~Geist, and D.~Peters,
		\textit{Optimal Bounds for the No-Show Paradox via SAT Solving},
		\textit{Mathematical Social Sciences}, 90 (2017), 18--27.
		
		\bibitem{fishburn1983}
		P.~C.~Fishburn and S.~J.~Brams,
		\textit{Paradoxes of Preferential Voting},
		Mathematics Magazine, 56(4) (1983), 207--214.
		
		\bibitem{moulin1988}
		H.~Moulin,
		\textit{Condorcet’s Principle Implies the No-Show Paradox},
		Journal of Economic Theory, 45(1) (1988), 53--64.
		
		\bibitem{rzazewski2014}
		K.~Rzążewski, W.~Słomczyński, and K.~Życzkowski,
		\textit{Każdy głos się liczy! Wędrówka przez krainę wyborów}, Warszawa, 2014.
		
	\end{thebibliography}
	
	\chapter*{Spis symboli i oznaczeń}
	
	\addcontentsline{toc}{chapter}{Spis symboli i oznaczeń}
	
	\begin{itemize}
		\item \(W\) — zbiór wszystkich wyborców, \(\#W = n\),
		\item \(K\) — zbiór wszystkich kandydatów, \(\#K = m\),

		\item \(\#A\) — liczność zbioru \(A\).
		
		\item \(\mathcal{P}(W)\) — zbiór potęgowy zbioru \(W\),
		\item \(\mathcal{E}(W) := \mathcal{P}(W) \setminus \{\emptyset\}\) — zbiór wszystkich elektoratów (niepustych podzbiorów zbioru \(W\)),
		
		\item \(\mathcal{R}\) — zbiór wszystkich relacji preferencji na zbiorze \(K\),
		\item \(\mathcal{R}^{V}\) — zbiór wszystkich profili preferencji określonych na elektoracie \(V\),
		
		\item \(\overset{w}{\succ}\) — relacja ścisłej preferencji wyborcy \(w \in W\),
		\item \(\overset{w}{\sim}\) — relacja równej preferencji wyborcy \(w \in W\),
		\item \(\overset{w}{\succcurlyeq}\) — relacja słabej preferencji wyborcy \(w \in W\),

		\Pytanie{Który ze znaków najlepiej opisuje słabą relację? 
			$$\succeq\qquad\succcurlyeq\qquad\succsim$$}
		
		\item \(R : V \to \mathcal{R}\) — profil preferencji na elektoracie \(V\),
		\item \(\Sigma = (V, R)\) — model wyborczy,
		
		\item \(M : \mathcal{R}^{V} \to \mathcal{P}(K)\) — metoda wyborcza,
		\item \(M : \mathcal{R}^{V} \to K\) — metoda rozstrzygająca (szczególny przypadek metody wyborczej),
		
		\item \(g_R : K \times K \to \mathbb{Z}\) — funkcja marginesu większości,
		
		\item \(T_R\) — ważony turniej indukowany przez profil preferencji \(R\),
		
		
	\end{itemize}
	
	\chapter*{Inne fragmenty do "wrzucenia" w prace}
	
	\begin{remark}
		Suma marginesów większościowych
		\[
		\sum_{x\neq a} g_R(a,x)
		\]
		indukuje ten sam porządek kandydatów co metoda Bordy.
	\end{remark}

	\Pytanie{Które z:
		\begin{itemize}
			\item Twierdzenie
			\item Definicja
			\item Lemat
			\item etc.
		\end{itemize}
	numerować 1.1, 1.2, 2.1 itd. a które np Tw. 1, Tw. 2, Tw. 3?
	}

	\Komentarz{Musze naprawić w dalszym ciągu wstawione "w złe miejsca" grafiki...}
	
	\Komentarz{Uczywam "na poniższej tabeli", a później jest ona podpisana jako "tablica". Czy mam to słownictwo uwspólnić?}
	
\end{document}